\subsection{Giới thiệu}

Phương pháp \tr{kernel} (kernel methods) được sử dụng rộng rãi trong học máy.
Đây là các kỹ thuật linh hoạt có thể mở rộng các thuật toán như SVM để xác
định các ranh giới quyết định phi tuyến. Các thuật toán khác chỉ phụ thuộc vào
tích trong giữa các điểm mẫu cũng có thể được mở rộng tương tự. Ý tưởng cốt
lõi dựa trên các hàm \tr{kernel}, vốn---dưới một số điều kiện kỹ thuật về
tính đối xứng và xác định dương---ngầm định nghĩa một tích trong trong không
gian nhiều chiều. Thay thế tích trong ban đầu trong không gian đầu vào bằng
\tr{kernel} xác định dương cho phép mở rộng các thuật toán như SVM thành phép
tách tuyến tính trong không gian đó, hay tương đương, thành phép tách phi tuyến
trong không gian đầu vào.

Cụ thể, ý tưởng là ánh xạ mỗi điểm dữ liệu $x \in \mathcal{X}$ vào một không
gian Hilbert $\mathbb{H}$ có số chiều cao hơn thông qua ánh xạ $\Phi \colon
\mathcal{X} \to \mathbb{H}$, nơi phép tách tuyến tính trở nên khả thi. Tuy
nhiên, $\mathcal{H}$ có thể có số chiều cực lớn---thậm chí vô hạn---khiến
việc tính tường minh $\Phi(x)$ là bất khả thi. Phương pháp \tr{kernel} giải quyết
điều này bằng cách thay thế trực tiếp tích trong $\langle \Phi(x), \Phi(x')
\rangle_{\mathbb{H}}$ bằng một hàm \tr{kernel} được tính hoàn toàn trong không
gian gốc $\mathcal{X}$, không phụ thuộc vào số chiều của $\mathbb{H}$.

\begin{defn}[\tr{Kernel}]
  Một hàm $K\colon \mathcal{X} \times \mathcal{X}$ được gọi là \tr{kernel}
  (kernel) trên $\mathcal{X}$.

  Ý tưởng là định nghĩa một \tr{kernel} $K$ sao cho với bất kỳ hai điểm
  $x, x'\in \mathcal{X}, K(x, x')$ bằng tích trong của các véc-tơ $\Phi(x)$
  và $\Phi(x')$:
  \begin{equation}
    \forall x, x' \in \mathcal{X}, \quad K(x, x') = \bigl\langle \Phi(x),
    \Phi(x') \bigr\rangle,
  \end{equation}
  với ánh xạ $\Phi(x)\colon \mathcal{X} \to \mathbb{H}$ vào một không gian
  Hilbert $\mathbb{H}$ được gọi là không gian đặc trưng. Vì tích trong là độ
  đo sự tương đồng của hai véc-tơ $K$ thường được hiểu là độ đo sự tương đồng
  giữa các phần các phần tử của không gian đầu vào $\mathcal{X}$.
\end{defn}\

Một ưu điểm quan trọng của \tr{kernel} $K$ là hiệu quả tính toán: $K$ thường
hiệu quả hơn đáng kể so với việc tính $\Phi(x)$ và tích trong $\mathbb{H}$.
Chúng ta sẽ thấy một số ví dụ phổ biến có thể đạt được trong đó việc tính
$K(x, x')$ có thể đạt được trong $O(N)$ trong khi tính
$\bigl\langle \Phi(x), \Phi(x') \bigr\rangle$ thường đòi hỏi
$O(\dim(\mathbb{H}))$ phép tính, với $\dim(\mathbb{H}) \gg N$. Hơn nữa,
trong một số trường hợp, số chiều của $\mathbb{H}$ là vô hạn.

Có lẽ lợi ích quan trọng hơn của hàm kernel $K$ là tính linh hoạt: không cần
định nghĩa hay tính toán tường minh ánh xạ $\Phi$. Kernel $K$ có thể được chọn
tùy ý miễn là sự tồn tại của $\Phi$ được đảm bảo, tức là $K$ thỏa mãn điều
kiện Mercer (xem định lý~\ref{thm:mercer-condition}).

\begin{thm}[Điều kiện Mercer]
  \label{thm:mercer-condition}
  Cho $\mathcal{X} \subset \mathbb{R^N}$ là một tập compact và $K \colon
  \mathcal{X} \times \mathcal{X} \to \mathbb{R}$ là một hàm liên tục và đối
  xứng. Khi đó, $K$ có thể khai triển hội tụ đều dưới dạng
  \begin{equation}
    K(x, x') = \sum_{n=0}^{\infty} a_n \phi_n(x)\phi_n(x'),
  \end{equation}
  với $a_n > 0$ khi và chỉ khi với mọi hàm bình phương khả tích $c \in
  L^2(\mathcal{X})$, điều kiện sau được thỏa mãn:
  \begin{equation}
    \iint_{\mathcal{X} \times \mathcal{X}} c(x)c(x')K(x,x')\,dx\,dx' \geq 0.
  \end{equation}
\end{thm}
Điều kiện này quan trọng để đảm bảo tính lồi của bài toán tối ưu trong các
thuật toán như SVM, từ đó đảm bảo hội tụ về cực tiểu toàn cục. Một điều kiện
tương đương với điều kiện Mercer dưới các giả thiết của định lý là kernel $K$
phải là \textit{\tr{kernel} đối xứng xác định dương}
(Positive Definite Symmetric, PDS). Tính chất này thực ra tổng quát hơn vì
đặc biệt nó không đòi hỏi bất kỳ giả thiết nào về $\mathcal{X}$.
